\section{Prover--verifier interaction}
\label{sec:interaction}

\definecolor{pvnavy}{HTML}{17365D}
\definecolor{pvtraffic}{HTML}{F3F4F6}
\definecolor{pvhint}{HTML}{A15C00}

\newcommand{\Nbadge}[1]{\textcolor{specblue}{\chN[#1]}}
\newcommand{\Hbadge}[1]{\textcolor{pvhint}{\chH[#1]}}
\newcommand{\Dbadge}[1]{\textcolor{specmuted}{\chD[#1]}}
\newcommand{\Splitbadge}[1]{%
  \textcolor{specblue}{\chN}\!/
  \textcolor{pvhint}{\chH}[#1]}

\newcommand{\idCS}{$\mathsf{CS}[h]$}
\newcommand{\idLINH}{$\mathsf{LIN}[h]$}
\newcommand{\idSUMIN}{$\mathsf{SC}_{\rm in}[H]$}
\newcommand{\idPESAT}{$\mathsf{PESAT}$}
\newcommand{\idSUMH}{$\mathsf{SUM}[H]$}
\newcommand{\idEVALH}{$\mathsf{EVAL}[H]$}
\newcommand{\idLINPACK}{$\mathsf{LIN}[f]$}
\newcommand{\idLINID}{$\mathsf{LIN}[\mathsf{id}]$}

\newcolumntype{L}{>{\raggedright\arraybackslash}p{2.04in}}
\newcolumntype{T}{>{\centering\arraybackslash\columncolor{pvtraffic}}p{0.65in}}
\newcolumntype{R}{>{\raggedright\arraybackslash}p{2.24in}}

\newenvironment{pvexchange}
  {\par\scriptsize
   \setlength{\tabcolsep}{3pt}%
   \setlength{\abovedisplayskip}{2pt}%
   \setlength{\belowdisplayskip}{2pt}%
   \renewcommand{\arraystretch}{1.02}%
   \begin{tabular}{@{}LTR@{}}
   \rowcolor{pvnavy}
   \multicolumn{1}{c}{\color{white}\bfseries Prover $\mathsf P$}
   &\multicolumn{1}{c}{\color{white}\bfseries Traffic}
   &\multicolumn{1}{c}{\color{white}\bfseries Verifier $\mathsf V$}
   \\
   \addlinespace[1pt]}
  {\end{tabular}\par}

\newenvironment{pvlongexchange}
  {\par\scriptsize
   \setlength{\LTleft}{0pt}\setlength{\LTright}{\fill}%
   \setlength{\LTpre}{0pt}\setlength{\LTpost}{0pt}%
   \setlength{\tabcolsep}{3pt}%
   \setlength{\abovedisplayskip}{2pt}%
   \setlength{\belowdisplayskip}{2pt}%
   \renewcommand{\arraystretch}{1.02}%
   \begin{longtable}{@{}LTR@{}}
   \rowcolor{pvnavy}
   \multicolumn{1}{c}{\color{white}\bfseries Prover $\mathsf P$}
   &\multicolumn{1}{c}{\color{white}\bfseries Traffic}
   &\multicolumn{1}{c}{\color{white}\bfseries Verifier $\mathsf V$}
   \\
   \addlinespace[1pt]
   \endfirsthead
   \rowcolor{pvnavy}
   \multicolumn{1}{c}{\color{white}\bfseries Prover $\mathsf P$}
   &\multicolumn{1}{c}{\color{white}\bfseries Traffic}
   &\multicolumn{1}{c}{\color{white}\bfseries Verifier $\mathsf V$}
   \\
   \addlinespace[1pt]
   \endhead}
  {\end{longtable}\par}

\newcommand{\Pmsg}[1]{\(\xrightarrow{\;\scriptstyle #1\;}\)}
\newcommand{\Vmsg}[1]{\(\xleftarrow{\;\scriptstyle #1\;}\)}
\newcommand{\Pturn}[1]{#1 & & \\[1pt]}
\newcommand{\Vturn}[1]{& & #1 \\[1pt]}
\newcommand{\Psend}[1]{& \Pmsg{#1} & \\[-1pt]}
\newcommand{\Vsend}[1]{& \Vmsg{#1} & \\[-1pt]}

\newenvironment{pvstep}[3][2.0in]
  {\Needspace{#1}%
   \par\addvspace{0.45\baselineskip}%
   \refstepcounter{clgstep}%
   \clgmarkstart
   \noindent
   {\normalsize\bfseries\sffamily\textcolor{clgtitlec}%
     {P\theclgstep/V\theclgstep.}~#2}%
   \nopagebreak\par\vspace{0.08\baselineskip}%
   #3%
   \clgboard%
   \nopagebreak}
  {\clgreserve\par\addvspace{0.12\baselineskip}}

\subsection{End-to-end proof decomposition and transcript channels}
\label{sec:transcript-channels}

The end-to-end proof consists of the configured PIOP proof followed by the
F2Z proof:
\begin{equation}
 \pi_{\rm e2e}:=\left(
   \pi_{\rm PIOP},
   \pi_{\rm F2Z}
 \right).
 \label{eq:end-to-end-proof-object}
\end{equation}
Verification of $\pi_{\rm PIOP}$ reduces the initial constraint statement to
the F2Z handoff
\begin{equation}
 \left\langle\bm v,\operatorname{vec}(H_q)\right\rangle_{\Fq}=\mu.
 \label{eq:piop-f2z-handoff}
\end{equation}
The pair $(\bm v,\mu)$ is produced by the PIOP verifier and is not repeated as
an F2Z proof field.

Within F2Z, write
\begin{equation}
 \pi_1:=\left(
   \pi_{1,{\rm in}},
   \bm\nu,\bm\mu
 \right),
 \qquad
 \pi_{1,{\rm in}}:=(g_i^{\rm shape})_{i=0}^{m_h-1}.
 \label{eq:pi1-proof-object}
\end{equation}
The semantic F2Z proof object is
\begin{equation}
 \pi_{\mathrm{F2Z}}
 =\left(
   \pi_1,
   \pi_{\rm GKR},\pi_{\rm MLE},
   \bm s,
   \pi_{\rm open}
 \right).
 \label{eq:proof-object}
\end{equation}
The fragment $\pi_{1,{\rm in}}$ contains the round polynomials of the standard
input Sumcheck that reduces Equation~\eqref{eq:piop-f2z-handoff} to
\begin{equation}
 \kappa\widetilde H_q(\bm s_h)=\tau.
 \label{eq:input-sumcheck-output}
\end{equation}
The point $\bm s_h$, scalars $\kappa$ and $\tau$, and factors
$(\bm a,\bm a')$ are verifier-derived and are not additional proof fields.
The ring-switch component is
$\bm s=(\bm s_\theta)_{\theta\in\Theta}$, with each
$\bm s_\theta\in\K^{128}$; the ordered set $\Theta$ and its coefficient
decomposition are public and are not proof fields.

Equations~\eqref{eq:end-to-end-proof-object}--\eqref{eq:proof-object} describe
semantic proof components.  The developer configuration fixes the encoding of
$\pi_{\rm PIOP}$, while Section~\ref{sec:wire-format} fixes the encoding of
$\pi_{\rm F2Z}$.  The F2Z v1 container partitions its transmitted records into
the NARG stream $N_{\rm F2Z}$ and the hint stream $H_{\rm F2Z}$:
\begin{equation}
 \pi_{\rm wire}:=(N_{\rm F2Z},H_{\rm F2Z}),\qquad
 \begin{aligned}
  \NargLen(\pi_{\rm F2Z})&:=|N_{\rm F2Z}|,&
  \HintLen(\pi_{\rm F2Z})&:=|H_{\rm F2Z}|,\\
  \PayloadLen(\pi_{\rm F2Z})&:=|N_{\rm F2Z}|+|H_{\rm F2Z}|,&
  \WireLen(\pi_{\rm F2Z})&:=|\HostEncode(N_{\rm F2Z},H_{\rm F2Z})|.
 \end{aligned}
 \label{eq:proof-channel-lengths}
\end{equation}

Equation~\eqref{eq:proof-channel-lengths} defines the two F2Z proof streams;
the following badges identify how each transmitted value is handled.

\begingroup\footnotesize
\renewcommand{\arraystretch}{1.05}
\begin{tabularx}{\linewidth}{@{}p{0.48in}p{0.90in}X@{}}
\toprule
\textbf{Badge}&\textbf{Source}&\textbf{Verifier action}\\
\midrule
$\Nbadge{x}$&NARG stream&Receive $x$ and verify it, and absorb it into the
verifier's transcript.\\
$\Hbadge{x}$&Hint stream&Receive $x$ and verify it, but do not absorb it into the
verifier's transcript.\\
$\Dbadge{x}$&No proof bytes&Receive nothing; $x$ is public or derived by the
verifier.\\
$\Splitbadge{\pi}$&Both streams&Receive each field from its specified stream.
Absorb NARG fields and verify Hint fields.\\
\bottomrule
\end{tabularx}
\endgroup

The verifier completes every required check and absorption before deriving the
next challenge.

\paragraph{Transcript absorption.}
Every absorbed value uses the encoding
\begin{equation}
 \Absorb\!\left(
   \mathsf{DOMAIN}
   \parallel\mathsf{PAYLOAD\_SIZE}
   \parallel\mathsf{PAYLOAD}
 \right).
 \label{eq:transcript-absorb}
\end{equation}
Section~\ref{sec:wire-format} specifies the canonical encoding of each domain
and payload, together with the transcript replay rules.  The domain identifies
the semantic meaning of the absorbed value.

Right ledgers contain post-step claims; \emph{internal} claims are shown only
when they explain why a later acceptance check cannot yet discharge the
original relation.

\setlength{\textwidth}{5.63in}
\setlength{\columnwidth}{5.63in}
\setlength{\linewidth}{5.63in}
\setlength{\hsize}{5.63in}
\setlength{\marginparwidth}{1.72in}
\setlength{\marginparsep}{0.10in}

% --------------------------------------------------------------------------
\begin{pvstep}{Bind the constraint statement and commitment}{%
  \clgnewrel{\idCS}
    {$(x;H_q)\in\mathcal R_{\rm CS}$}%
  \clgdef{$\rho$}{$\rho=\Root(O)$}%
}
\begin{pvexchange}
\Pturn{Construct the committed and PIOP-facing witness views:
\(\begin{aligned}
 \pi_2(\bm h)&=\bm M\pi_2(\bm f),\\
 \operatorname{vec}(H)&=\pi_2(\bm h),\\
 \operatorname{vec}(F)&=\pi_2(\bm f),\qquad P:=\Pack(F),\\
 O&:=\Enc_{\mathcal C}(P),\qquad \rho:=\Root(O).
\end{aligned}\)
and retain $(\bm h,H,F,P,O)$.}
\Psend{\Dbadge{x,\rho}}
\Vturn{Validate the constraint statement, the developer configuration, and
the commitment root.  Bind $(x,\rho)$ as required by the configured PIOP.}
\end{pvexchange}
\end{pvstep}

\smallskip
\noindent\textbf{Configured matrix indexing.}
The developer configuration fixes a bijection between each Boolean address
$z\in\bitsset^{m_h}$ and a pair
\[
 z=(r,c)\in\bitsset^{d_r}\times\bitsset^{d_c},
 \qquad m_h=d_r+d_c.
\]
Using this fixed address split, define
\[
 H[r,c]:=\bigl[\pi_2(\bm h)\bigr]_z,
 \qquad
 \operatorname{vec}(H)=\pi_2(\bm h).
\]
Thus $H$ is a matrix view of the same witness bits, not a new witness.  The
tables $H_{\mathbb Z}$, $H_q$, and $H_{\K}$ contain these same entries embedded
into $\mathbb Z$, $\Fq$, and $\K$, respectively.  We use brackets for table
entries, such as $H_q[r,c]$, and parentheses only for multilinear-extension
evaluations, such as $\widetilde H_q(\bm x_r,\bm x_c)$.

% --------------------------------------------------------------------------
\begin{pvstep}[2.65in]{Run the PIOP}{%
  \clgedge[PIOP]{$\Pi_{\rm PIOP}$}
    {$\mathcal R_{\rm CS}$}
    {$\mathcal R_{\rm LinBitsDual}[\bm h]$}%
  \clgdone{\idCS}%
  \clgnewrel{\idLINH}
    {$\langle\bm v,\operatorname{vec}(H_q)\rangle_{\Fq}=\mu$}%
}
\begin{pvexchange}
\Pturn{Run the configured PIOP on the constraint statement and witness:
\[
 \pi_{\rm PIOP}\leftarrow
 \Pi_{\rm PIOP}.\mathsf{Prove}
   (\mathsf{pp}_{\rm PIOP},x,\operatorname{vec}(H_q)).
\]}
\Psend{\pi_{\rm PIOP}}
\Vturn{Run the matching PIOP verifier:
\[
 (\mathsf{ok}_{\rm PIOP},\bm v,\mu)\leftarrow
 \Pi_{\rm PIOP}.\mathsf{Verify}
   (\mathsf{pp}_{\rm PIOP},x,\pi_{\rm PIOP}).
\]
Require $\mathsf{ok}_{\rm PIOP}=1$.  Its output is the F2Z handoff
\[
 \boxed{\left\langle\bm v,\operatorname{vec}(H_q)\right\rangle_{\Fq}=\mu}.
\]}
\Vsend{\Dbadge{\bm v,\mu}}
\end{pvexchange}
\end{pvstep}

% --------------------------------------------------------------------------
\begin{pvstep}[3.05in]{Reduce the PIOP handoff with Sumcheck}{%
  \clgedge[Sumcheck]{$\Pi_{1,{\rm in}}$}
    {$\mathcal R_{\rm LinBitsDual}[\bm h]$}
    {$\kappa\widetilde H_q(\bm s_h)=\tau$}%
  \clgdone{\idLINH}%
  \clgnewinternal{\idSUMIN}
    {$\kappa\widetilde H_q(\bm s_h)=\tau$}%
}
\begin{pvexchange}
\Vturn{Initialize the F2Z transcript by binding its domain, handoff statement,
configuration, and commitment root.}
\Pturn{Run the standard $m_h$-round Sumcheck on the Lin-BitsDual claim from
Section~\ref{sec:linbits-instance}:
\[
 \begin{aligned}
  \bm X&=(\bm X_r,\bm X_c),\\
  G_{\rm shape}(\bm X)
    &:=\widetilde{\bm v}(\bm X_r,\bm X_c)\\[-2pt]
    &\phantom{:=}{}\cdot\widetilde H_q(\bm X_r,\bm X_c),\\
  \mu&=\sum_{r,c}G_{\rm shape}(r,c),\\
  \pi_{1,{\rm in}}&:=(g_i^{\rm shape})_{i=0}^{m_h-1}.
 \end{aligned}
\]}
\Psend{\Nbadge{\pi_{1,{\rm in}}}}
\Vturn{Verify the Sumcheck transcript.  It yields a terminal point
$\bm s_h=(\bm s_r,\bm s_c)$ and value $\tau$ satisfying
\[
 \tau=\widetilde{\bm v}(\bm s_r,\bm s_c)
      \widetilde H_q(\bm s_r,\bm s_c),
 \qquad
 \kappa:=\widetilde{\bm v}(\bm s_r,\bm s_c),
\]
and derive
\(\begin{aligned}
 a_r&:=\eq_{\Fq}(r,\bm s_r),\\
 a'_c&:=\kappa\eq_{\Fq}(c,\bm s_c).
\end{aligned}\)
The resulting terminal claim is
\[
 \boxed{
  \kappa\widetilde H_q(\bm s_r,\bm s_c)
  =\tau
  =\sum_c a'_c\sum_r a_rH_q[r,c]
 }.
\]
Absorb the derived seam and derive the canonical integer lifts
$\gamma_r:=\operatorname{lift}_q(a_r)$.}
\Vsend{\Dbadge{\bm s_h,\kappa,\tau}}
\end{pvexchange}
\end{pvstep}

% --------------------------------------------------------------------------
\begin{pvstep}[3.15in]{Fold each column and encode the folds as product claims}{%
  \clgdone{\idSUMIN}%
  \clgnewrel{\idPESAT}{$\mathcal R_{\rm PESAT}[\bm h]$}%
  \clgnote{\idPESAT}{%
    $\nu_c=\prod_r
      \bigl(1+(g^{\gamma_r}-1)H_{\K}[r,c]\bigr)$}%
  \clgdef{$P$}{$P=\Pack(F)$}%
}
\begin{pvlongexchange}
\Pturn{For every column $c$, compute
\(\begin{aligned}
 \widehat\mu_c&:=\sum_r\gamma_rH_{\mathbb Z}[r,c],\\
 \mu_c&:=\widehat\mu_c,\\
 \nu_c&:=g^{\mu_c}\in\K^\times.
\end{aligned}\)}
\Psend{\Nbadge{\bm\nu}}
\Psend{\Hbadge{\bm\mu}}
\Vturn{Before deriving any dependent challenge, require
\[
 |\bm\mu|=2^{d_c},
\]
and, for every $c$,
\[
 0\le\mu_c<2^{127},\qquad g^{\mu_c}=\nu_c.
\]
Then close the retained Sumcheck claim by checking
\[
 \boxed{\sum_c a'_c\pi_q(\mu_c)=\tau}.
\]}
\Vturn{After accepting the hint, derive
\(\begin{aligned}
 \zeta&\leftarrow\Squeeze_{\K}^{d_c}(\tr),\\
 e_0&:=\MLE{\bm\nu}(\zeta).
\end{aligned}\)}
\Vsend{\Dbadge{\zeta}}
\Vturn{Record the successor claim for P5/V5:
\(\begin{aligned}
 Q_{r,c}&:=1+(g^{\gamma_r}-1)H_{\K}[r,c],\\
 \nu_c&=\prod_r Q_{r,c}.
\end{aligned}\)
This equality is not checked in P4/V4, and no entry of $H_{\K}$ is sent.}
\end{pvlongexchange}
\end{pvstep}

% --------------------------------------------------------------------------
\begin{pvstep}[2.35in]{Run the GKR reduction}{%
  \clgedge[GKR subprotocol]{$\Pi_{\rm GKR}$}
    {$\mathcal R_{\rm PESAT}[\bm h]$}
    {$\mathcal R_{\rm sum}^{h}$}%
  \clgdone{\idPESAT}%
  \clgnewrel{\idSUMH}
    {$S(H_{\K})=\sigma$}%
  \clgnote{\idSUMH}{%
    $S(H_{\K})=\sum_z\omega_{\xi,z}H_{\K}[z]$}%
}
\begin{pvexchange}
\Pturn{Set
\(\displaystyle
 \mathcal I_{\rm GKR}:=
 \mathsf{GKRInput}(\bm\gamma,\bm\nu,\zeta,e_0)
\)
and compute
\(\displaystyle
 \pi_{\rm GKR}\leftarrow
 \mathsf{GKR.Prove}
 (\mathsf{pp}_{\rm GKR},\mathcal I_{\rm GKR},H_{\K};\tr)\).}
\Psend{\Splitbadge{\pi_{\rm GKR}}}
\Vturn{Compute
\(\begin{aligned}
 (\tr',\bm\omega_\xi,\sigma)&\leftarrow
   \mathsf{GKR.Verify}\\
 &\quad(\mathsf{pp}_{\rm GKR},\mathcal I_{\rm GKR},
        \pi_{\rm GKR};\tr),
\end{aligned}\)
require success and the exact subordinate endpoint, then set
$\tr\leftarrow\tr'$.  The successor is
\(\displaystyle
 \sum_{z\in\bitsset^{m_h}}
   \omega_{\xi,z}H_{\K}[z]=\sigma.
\)}
\end{pvexchange}
GKR operates on $H_{\K}$ and does not query $O$.
\end{pvstep}

% --------------------------------------------------------------------------
\begin{pvstep}[3.05in]{Reduce the GKR sum claim to an MLE claim}{%
  \clgedge[row LinCheck]{$\Pi_{\rm MLE}$}
    {$\mathcal R_{\rm sum}^{h}$}
    {$\mathcal R_{\rm eval}^{h}$}%
  \clgdone{\idSUMH}%
  \clgnewrel{\idEVALH}
    {$\widetilde H_{\K}(\bm r_h)=\eta$}%
  \clgundef{$P$}%
}
\begin{pvexchange}
\Pturn{Run the standard $m_h$-round Sumcheck on the weighted-sum claim:
\[
 \begin{aligned}
  G_{\rm MLE}(\bm X)
    &:=\widetilde{\bm\omega_\xi}(\bm X)
       \widetilde H_{\K}(\bm X),\\
  \sigma&=\sum_{z\in\bitsset^{m_h}}G_{\rm MLE}(z),\\
  \pi_{\rm MLE}
    &:=\bigl((g_i^{\rm MLE})_{i=0}^{m_h-1},\eta\bigr),
 \end{aligned}
\]
where the prover supplies
$\eta:=\widetilde H_{\K}(\bm r_h)$ at the transcript-derived terminal point
$\bm r_h$.}
\Psend{\Splitbadge{\pi_{\rm MLE}}}
\Vturn{Verify the Sumcheck transcript.  It yields a terminal point
$\bm r_h$ and value $\theta$.  Compute
\[
 \lambda:=\widetilde{\bm\omega_\xi}(\bm r_h)
\]
and require the terminal check
\[
 \boxed{\theta=\lambda\eta}.
\]
The resulting successor claim is
\[
 \boxed{\widetilde H_{\K}(\bm r_h)=\eta}.
\]}
\end{pvexchange}
The row LinCheck operates on $H_{\K}$ and does not query $O$.
\end{pvstep}

% --------------------------------------------------------------------------
\begin{pvstep}[2.25in]{Rewrite $\langle\bm a,\bm M\bm f\rangle$ as
$\langle\bm M^{\mathsf T}\bm a,\bm f\rangle$}{%
  \clgedge[adjoint identity]
    {$\begin{gathered}
       \langle\bm a,\bm M_{\K}\bm f_{\K}\rangle_{\K}\\[-2pt]
       =\langle\bm M_{\K}^{\mathsf T}\bm a,\bm f_{\K}\rangle_{\K}
     \end{gathered}$}
    {$\mathcal R_{\rm eval}^{h}$}
    {$\LIN(\K,\mathcal C,\Pack)$}%
  \clgdone{\idEVALH}%
  \clgnewrel{\idLINPACK}{$\langle\bm a_f,\bm f_{\K}\rangle_{\K}=\eta$}%
}
\begin{pvlongexchange}
\multicolumn{3}{@{}p{5.05in}@{}}{
Both parties set $\bm a[z]:=\eq(z,\bm r_h)$ and use
$H_{\K}=\bm M_{\K}\bm f_{\K}$ to rewrite
\[
 \eta
 =\langle\bm a,H_{\K}\rangle_{\K}
 =\langle\bm a,\bm M_{\K}\bm f_{\K}\rangle_{\K}
 =\langle\bm M_{\K}^{\mathsf T}\bm a,\bm f_{\K}\rangle_{\K}.
\]
Set $\bm a_f:=\bm M_{\K}^{\mathsf T}\bm a$.  No message or challenge is
needed.  In the configured packing order, both parties also view this same
coefficient as
\[
 A_f[v,j]:=\bm a_f[128j+v].
\]
This is only a change of coordinates; the claim remains
$\langle\bm a_f,\bm f_{\K}\rangle_{\K}=\eta$.}
\end{pvlongexchange}
\end{pvstep}

% --------------------------------------------------------------------------
\begin{pvstep}[4.35in]{Ring-switch the adjoint claim into the committed packing}{%
  \clgedge[$\Pi_{\rm iopp}$ internal]{$\Pi_{\rm ring}$}
    {$\LIN(\K,\mathcal C,\Pack)$}
    {$\LIN(\K,\mathcal C,\mathsf{id})$}%
  \clgdone{\idLINPACK}%
  \clgnewrel{\idLINID}{$\langle B,P\rangle_{\K}=\beta$}%
  \clgdef{$P$}{$P=\Pack(F)$}%
}
\begin{pvlongexchange}
\multicolumn{3}{@{}p{5.05in}@{}}{
From the public coefficient $A_f$, both parties derive the canonical ordered
decomposition
\[
 A_f[v,j]=\sum_{\theta\in\Theta}A_\theta[v]B_\theta[j]
\]
from Section~\ref{sec:ring-switch-coefficients}.  This represents the same
inner product produced by P7/V7; it is not another reduction or proof claim.}
\\[2pt]
\Pturn{For every $\theta\in\Theta$ and $0\le v<128$, compute
\[
 s_{\theta,v}:=\sum_{j<2^{m_p}}B_\theta[j]F_v[j].
\]}
\Psend{\Nbadge{\bm s_\theta}}
\Vturn{Require $|\bm s_\theta|=128$ for every $\theta$ and anchor all
messages to the adjoint claim:
\[
 \boxed{
  \sum_{\theta\in\Theta}\sum_{v=0}^{127}
    A_\theta[v]s_{\theta,v}=\eta
 }.
\]}
\Vturn{After accepting every $\bm s_\theta$, derive
\(\displaystyle
 r''\leftarrow\Squeeze_{\K}^{7}(\tr).
\)}
\Vsend{\Dbadge{r''}}
\multicolumn{3}{@{}>{\raggedright\arraybackslash}p{5.05in}@{}}{
With $[x]_u$ denoting coordinate $u$ in the fixed $\K/\Ftwo$ basis, both
parties derive, for every $\theta$,
\(\begin{aligned}
 t_{\theta,u}&:=\sum_v[s_{\theta,v}]_u\beta_v,\\
 \Phi_{r''}(x)&:=\sum_u\eq(u,r'')[x]_u,\\
 \beta_\theta&:=\sum_u\eq(u,r'')t_{\theta,u},\\
 \omega_\theta[j]&:=\Phi_{r''}(B_\theta[j]),
\end{aligned}\)
which gives $\langle\bm\omega_\theta,P\rangle_{\K}=\beta_\theta$.}
\\[2pt]
\Vturn{For $\theta\in\Theta$ in order, derive
\(\displaystyle
 \lambda_\theta\leftarrow\Squeeze_{\K}(\tr).
\)}
\Vsend{\Dbadge{(\lambda_\theta)_{\theta\in\Theta}}}
\Pturn{Set
\(\displaystyle
 B:=\sum_{\theta\in\Theta}\lambda_\theta\bm\omega_\theta,
 \qquad
 \beta:=\sum_{\theta\in\Theta}\lambda_\theta\beta_\theta.
\)}
\Vturn{Independently derive the same $(B,\beta)$.}
\end{pvlongexchange}
\end{pvstep}

% --------------------------------------------------------------------------
\begin{pvstep}[2.65in]{Open the packed claim}{%
  \clgedge[$\Pi_{\rm iopp}$ internal]{$\Pi_{\rm open}$}
    {$\LIN(\K,\mathcal C,\mathsf{id})$}
    {$\mathcal R_{\rm triv}$}%
  \clgdone{\idLINID}%
  \clgundef{$P$}%
  \clgundef{$\rho$}%
}
\begin{pvlongexchange}
\Pturn{Absorb the opening target $(m_p,\beta)$.}
\Vturn{Retain the locally derived $B$ from P8/V8 and absorb the same opening
target.}
\Pturn{Compute
\(\displaystyle
 \pi_{\rm open}\leftarrow
 \mathsf{RecProve}(\mathsf{pp},\rho,B,\beta,P;\tr).
\)}
\Psend{\Splitbadge{\pi_{\rm open}}}
\Vturn{Require
\(\begin{aligned}
 &\mathsf{RecVerify}
   (\mathsf{pp},\rho,B,\beta;\tr,c_N,c_H)=1,\\
 &\mathsf{NoPendingHint}=1,\\
 &\mathsf{NargEOF}(c_N)=\mathsf{HintEOF}(c_H)=1,\\
 &\mathsf{HostEOF}(\pi_{\rm host})=1,\\
 &\mathsf{Counts}(c_N,c_H)=(n_N,n_H).
\end{aligned}\)}
\end{pvlongexchange}
\end{pvstep}

\setlength{\textwidth}{7.45in}
\setlength{\columnwidth}{7.45in}
\setlength{\linewidth}{7.45in}
\setlength{\hsize}{7.45in}

P9/V9 follows the recursive-opening ledger in
Section~\ref{sec:wire-recursive-ledger}.  Its proof spans both streams; the
whole $\pi_{\rm open}$ \mustnot{} be placed on the hint channel.
